\documentclass[../DoAn.tex]{subfiles}
\begin{document}

% =====================================================================
% CHƯƠNG 9 — KẾT LUẬN VÀ HƯỚNG PHÁT TRIỂN
% =====================================================================

\section{Kết luận}
% Đối chiếu kết quả đạt được với 3 mục tiêu ở Chương 1.

Trong đồ án này, chúng tôi nghiên cứu bài toán trích xuất điện tâm đồ thai nhi không xâm lấn (NI-fECG) từ tín hiệu ECG bụng của mẹ (maECG) trong điều kiện SNR thấp và có chồng lấp mạnh giữa mECG–fECG. Dựa trên nền tảng W-NETR (hai nhánh UNETR 1D và cơ chế loại bỏ đặc trưng mECG trong latent space), mô hình đề xuất \textbf{KAN-WNETR} được xây dựng bằng cách \textbf{thay thế khối MLP/FFN trong Transformer encoder bằng KAN gốc} (learnable activation theo spline), trong khi vẫn giữ nguyên cơ chế self-attention và kiến trúc chữ W.

Thực nghiệm trên bộ dữ liệu tổng hợp FECGSYNDB với nhiều ngữ cảnh sinh lý/nhiễu (C0--C5) và các mức SNR khác nhau cho thấy mô hình có khả năng tái tạo tín hiệu fECG ổn định. Về chất lượng dạng sóng, KAN-WNETR đạt \textbf{SSIM tổng 0.966}, nhỉnh hơn kết quả báo cáo của W-NETR (0.960), đồng thời thể hiện mức cải thiện rõ ở các kịch bản nhiễu phức tạp. Về phát hiện đỉnh R của fECG, mô hình đạt \textbf{Recall 97.53\%}, \textbf{Precision 97.94\%} và \textbf{F1 97.73\%}, phản ánh sự cân bằng tốt giữa giảm bỏ sót và giảm báo sai trong môi trường nhiễu.

% TODO: bổ sung 1-2 câu đối chiếu trực tiếp với 3 mục tiêu ở Chương 1 (gạch đầu dòng):
\begin{itemize}
  \item Mục tiêu 1 (mô hình KAN-WNETR): đã đạt — nêu kết quả SSIM/F1.
  \item Mục tiêu 2 (triển khai biên thời gian thực): đã đạt — INT8 trên Pi 5, $\sim$8--12$\times$ thời gian thực.
  \item Mục tiêu 3 (hệ thống theo dõi): đã đạt — thiết bị + server + dashboard lâm sàng (Chương~\ref{chapter:System}).
\end{itemize}

% ---------------------------------------------------------------------
\section{Hạn chế}
% TODO: viết thành văn xuôi (giữ nội dung cũ, gạch đầu dòng định hướng)
\begin{itemize}
  \item Kết quả chủ yếu kiểm chứng trên dữ liệu mô phỏng (chưa có dữ liệu lâm sàng thực).
  \item Hiệu năng phát hiện fQRS còn khoảng cách so với trần SOTA (W-NETR) ở một số thiết lập.
  \item Chất lượng tín hiệu tụt sau lượng tử hoá INT8 so với FP32.
  \item Ablation MLP-FFN vs KAN-FFN cần được hoàn tất để định lượng riêng đóng góp của KAN.
\end{itemize}

% ---------------------------------------------------------------------
\section{Hướng phát triển}
% TODO: viết thành văn xuôi (gạch đầu dòng định hướng)
\begin{itemize}
  \item Thu \textbf{dữ liệu lâm sàng thực} (chính bộ thu ADS1293 của đề tài có thể đảm nhiệm).
  \item Mở rộng đa kênh; ablation sâu hơn; quantization-aware training (QAT) để thu hẹp khoảng cách FP32$\leftrightarrow$INT8.
  \item Tối ưu chiến lược huấn luyện/loss và hậu xử lý peak-picking.
  \item Tiến tới chuẩn thiết bị y tế.
\end{itemize}

% ---------------------------------------------------------------------
\section{Bài học kinh nghiệm}
% TODO: viết thành văn xuôi (gạch đầu dòng định hướng)
\begin{itemize}
  \item Kinh nghiệm về quy trình nghiên cứu, thực nghiệm có đối chứng và triển khai biên.
\end{itemize}

\end{document}
